\QHeading{第1题　口算。}
\par\noindent


\begin{oralgrid}
\OralItem{1}{30\times160}{4800}
\OralItem{2}{280\div20}{14}
\OralItem{3}{450+560}{1010}
\OralItem{4}{720-380}{340}
\OralItem{5}{125\times40}{5000}
\OralItem{6}{630\div70}{9}
\OralItem{7}{15\times8}{120}
\OralItem{8}{400\div50}{8}
\OralItem{9}{3.27+5.63}{8.9}
\OralItem{10}{9.85-4.65}{5.2}
\OralItem{11}{0.7+0.33}{1.03}
\OralItem{12}{10-2.81}{7.19}\end{oralgrid}
\par

\QHeading{口算思路补充（对应上表题号）}

\Answer{上面列出得数，下面按小题给出不同计算思路。计数单位法适用于小数，整数题可比较凑整、拆分或商不变等方法。}
\MethodTitle{第2小题：$280\div 20$}

\MethodBlock{口算思路一：同时去掉相同个数的0}{}{\begin{aligned}[t]&\phantom{={}}280\div 20\\
&=(280\div 10)\div (20\div 10)\\
&=28\div 2\\
&=14\end{aligned}}

\MethodBlock{口算思路二：拆分被除数}{}{\begin{aligned}[t]&\phantom{={}}280\div 20\\
&=(200+80)\div 20\\
&=200\div 20+80\div 20\\
&=10+4\\
&=14\end{aligned}}
\MethodTitle{第4小题：$720-380$}

\MethodBlock{口算思路一：补差法}{从减数先补到一个整数，再补到被减数，两段差相加。}{\begin{aligned}[t]&\phantom{={}}720-380\\
&=(390-380)+(720-390)\\
&=10+330\\
&=340\end{aligned}}
\MethodTitle{第5小题：$125\times 40$}

\MethodBlock{口算思路一：拆因数，凑整计算}{把一个因数拆成两个因数，先凑成整百或整千。}{\begin{aligned}[t]&\phantom{={}}125\times 40\\
&=125\times (8\times 5)\\
&=(125\times 8)\times 5\\
&=1000\times 5\\
&=5000\end{aligned}}

\MethodBlock{口算思路二：按数位拆分计算}{把一个因数拆成整十、整百数与余下部分之和。}{\begin{aligned}[t]&\phantom{={}}125\times 40\\
&=40\times (100+25)\\
&=40\times 100+40\times 25\\
&=4000+1000\\
&=5000\end{aligned}}
\MethodTitle{第6小题：$630\div 70$}

\MethodBlock{口算思路一：同时去掉相同个数的0}{}{\begin{aligned}[t]&\phantom{={}}630\div 70\\
&=(630\div 10)\div (70\div 10)\\
&=63\div 7\\
&=9\end{aligned}}

\MethodBlock{口算思路二：拆分除数，连续相除}{}{\begin{aligned}[t]&\phantom{={}}630\div 70\\
&=630\div (2\times 35)\\
&=630\div 2\div 35\\
&=315\div 35\\
&=9\end{aligned}}
\MethodTitle{第7小题：$15\times 8$}

\MethodBlock{口算思路一：凑整后补偿}{把接近整十、整百的因数凑整，再补回或减去多算的部分。}{\begin{aligned}[t]&\phantom{={}}15\times 8\\
&=15\times (10-2)\\
&=15\times 10-15\times 2\\
&=150-30\\
&=120\end{aligned}}

\MethodBlock{口算思路二：按数位拆分计算}{把一个因数拆成整十、整百数与余下部分之和。}{\begin{aligned}[t]&\phantom{={}}15\times 8\\
&=8\times (10+5)\\
&=8\times 10+8\times 5\\
&=80+40\\
&=120\end{aligned}}
\MethodTitle{第8小题：$400\div 50$}

\MethodBlock{口算思路一：同时去掉相同个数的0}{}{\begin{aligned}[t]&\phantom{={}}400\div 50\\
&=(400\div 10)\div (50\div 10)\\
&=40\div 5\\
&=8\end{aligned}}

\MethodBlock{口算思路二：利用商不变的性质}{被除数和除数同时乘同一个非零数，商不变。}{\begin{aligned}[t]&\phantom{={}}400\div 50\\
&=(400\times 2)\div (50\times 2)\\
&=800\div 100\\
&=8\end{aligned}}
\MethodTitle{第9小题：$3.27+5.63$}

\MethodBlock{口算思路一：按相同计数单位计算}{先看成若干个百分之一，按整数加减，最后还原成小数。}{\begin{aligned}[t]&\phantom{={}}3.27+5.63\\
&=(327+563)\div 100\\
&=890\div 100\\
&=8.9\end{aligned}}
\MethodTitle{第10小题：$9.85-4.65$}

\MethodBlock{口算思路一：按相同计数单位计算}{先看成若干个百分之一，按整数加减，最后还原成小数。}{\begin{aligned}[t]&\phantom{={}}9.85-4.65\\
&=(985-465)\div 100\\
&=520\div 100\\
&=5.2\end{aligned}}

\MethodBlock{口算思路二：补差法}{从减数先补到一个整数，再补到被减数，两段差相加。}{\begin{aligned}[t]&\phantom{={}}9.85-4.65\\
&=(5-4.65)+(9.85-5)\\
&=0.35+4.85\\
&=5.2\end{aligned}}
\MethodTitle{第11小题：$0.7+0.33$}

\MethodBlock{口算思路一：按相同计数单位计算}{先看成若干个百分之一，按整数加减，最后还原成小数。}{\begin{aligned}[t]&\phantom{={}}0.7+0.33\\
&=(70+33)\div 100\\
&=103\div 100\\
&=1.03\end{aligned}}
\MethodTitle{第12小题：$10-2.81$}

\MethodBlock{口算思路一：按相同计数单位计算}{先看成若干个百分之一，按整数加减，最后还原成小数。}{\begin{aligned}[t]&\phantom{={}}10-2.81\\
&=(1000-281)\div 100\\
&=719\div 100\\
&=7.19\end{aligned}}

\MethodBlock{口算思路二：补差法}{从减数先补到一个整数，再补到被减数，两段差相加。}{\begin{aligned}[t]&\phantom{={}}10-2.81\\
&=(3-2.81)+(10-3)\\
&=0.19+7\\
&=7.19\end{aligned}}


\par\noindent

\QHeading{第2题　列竖式计算。}

\Problem{（1）$35.68+14.32=50$}
\par


\MethodTitle{方法一：按题意列竖式}
\Answer{\begin{tikzpicture}[x=1em,y=1em,baseline=(current bounding box.center),draw=answerred,text=answerred]
\node[anchor=base,inner sep=0pt] at (0.000,0) {$3$};
\node[anchor=base,inner sep=0pt] at (0.800,0) {$5$};
\node[anchor=base,inner sep=0pt] at (1.600,0) {$.$};
\node[anchor=base,inner sep=0pt] at (1.920,0) {$6$};
\node[anchor=base,inner sep=0pt] at (2.720,0) {$8$};
\node[anchor=base,inner sep=0pt] at (0.000,-1.1) {$1$};
\node[anchor=base,inner sep=0pt] at (0.800,-1.1) {$4$};
\node[anchor=base,inner sep=0pt] at (1.600,-1.1) {$.$};
\node[anchor=base,inner sep=0pt] at (1.920,-1.1) {$3$};
\node[anchor=base,inner sep=0pt] at (2.720,-1.1) {$2$};
\node[anchor=base,inner sep=0pt] at (0.000,-2.3) {$5$};
\node[anchor=base,inner sep=0pt] at (0.800,-2.3) {$0$};
\node[anchor=base,inner sep=0pt] at (1.600,-2.3) {$.$};
\node[anchor=base,inner sep=0pt] at (1.920,-2.3) {$0$};
\node[anchor=base,inner sep=0pt] at (2.720,-2.3) {$0$};
\node[anchor=base,inner sep=0pt] at (-.8,-1.1) {$+$};
\draw[line width=.5pt] (-1.1,-1.52)--(3.240,-1.52);
\node[anchor=base,inner sep=0pt,font=\normalsize] at (1.920,.77) {$1$};
\node[anchor=base,inner sep=0pt,font=\normalsize] at (0.800,.77) {$1$};
\node[anchor=base,inner sep=0pt,font=\normalsize] at (0.000,.77) {$1$};
\end{tikzpicture}}

\MethodBlock{方法二：拆成整数部分和小数部分}{上方竖式及题目要求的验算仍保留；本法用于核对。}{\begin{aligned}[t]&\phantom{={}}35.68+14.32\\
&=35.68+14+0.32\\
&=49.68+0.32\\
&=50\end{aligned}}

\MethodBlock{方法三：按相同计数单位计算}{上方竖式及题目要求的验算仍保留；本法用于核对。先看成若干个百分之一，按整数加减，最后还原成小数。}{\begin{aligned}[t]&\phantom{={}}35.68+14.32\\
&=(3568+1432)\div 100\\
&=5000\div 100\\
&=50\end{aligned}}
\Problem{（2）$100-46.75=53.25$}

\MethodTitle{方法一：按题意列竖式}
\Answer{\begin{tikzpicture}[x=1em,y=1em,baseline=(current bounding box.center),draw=answerred,text=answerred]
\node[anchor=base,inner sep=0pt] at (0.000,0) {$1$};
\node[anchor=base,inner sep=0pt] at (0.800,0) {$0$};
\node[anchor=base,inner sep=0pt] at (1.600,0) {$0$};
\node[anchor=base,inner sep=0pt] at (2.400,0) {$.$};
\node[anchor=base,inner sep=0pt] at (2.720,0) {$0$};
\node[anchor=base,inner sep=0pt] at (3.520,0) {$0$};
\node[anchor=base,inner sep=0pt] at (0.800,-1.1) {$4$};
\node[anchor=base,inner sep=0pt] at (1.600,-1.1) {$6$};
\node[anchor=base,inner sep=0pt] at (2.400,-1.1) {$.$};
\node[anchor=base,inner sep=0pt] at (2.720,-1.1) {$7$};
\node[anchor=base,inner sep=0pt] at (3.520,-1.1) {$5$};
\node[anchor=base,inner sep=0pt] at (0.800,-2.3) {$5$};
\node[anchor=base,inner sep=0pt] at (1.600,-2.3) {$3$};
\node[anchor=base,inner sep=0pt] at (2.400,-2.3) {$.$};
\node[anchor=base,inner sep=0pt] at (2.720,-2.3) {$2$};
\node[anchor=base,inner sep=0pt] at (3.520,-2.3) {$5$};
\node[anchor=base,inner sep=0pt] at (-.8,-1.1) {$-$};
\draw[line width=.5pt] (-1.1,-1.52)--(4.040,-1.52);
\draw[line width=.35pt] (-0.190,.03)--(0.190,.47);
\node[anchor=base,inner sep=0pt,font=\normalsize] at (0.000,.77) {$0$};
\draw[line width=.35pt] (0.610,.03)--(0.990,.47);
\node[anchor=base,inner sep=0pt,font=\normalsize] at (0.800,.77) {$9$};
\draw[line width=.35pt] (1.410,.03)--(1.790,.47);
\node[anchor=base,inner sep=0pt,font=\normalsize] at (1.600,.77) {$9$};
\draw[line width=.35pt] (2.530,.03)--(2.910,.47);
\node[anchor=base,inner sep=0pt,font=\normalsize] at (2.720,.77) {$9$};
\draw[line width=.35pt] (3.330,.03)--(3.710,.47);
\node[anchor=base,inner sep=0pt,font=\normalsize] at (3.520,.77) {$10$};
\end{tikzpicture}}

\MethodBlock{方法二：多减再加回}{上方竖式及题目要求的验算仍保留；本法用于核对。减去凑成的整数后，把多减的部分加回来。}{\begin{aligned}[t]&\phantom{={}}100-46.75\\
&=100-47+0.25\\
&=53+0.25\\
&=53.25\end{aligned}}

\MethodBlock{方法三：拆成整数部分和小数部分}{上方竖式及题目要求的验算仍保留；本法用于核对。}{\begin{aligned}[t]&\phantom{={}}100-46.75\\
&=100-46-0.75\\
&=54-0.75\\
&=53.25\end{aligned}}
\Problem{（3）$28.7+54.36=83.06$}

\MethodTitle{方法一：按题意列竖式}
\Answer{\begin{tikzpicture}[x=1em,y=1em,baseline=(current bounding box.center),draw=answerred,text=answerred]
\node[anchor=base,inner sep=0pt] at (0.000,0) {$2$};
\node[anchor=base,inner sep=0pt] at (0.800,0) {$8$};
\node[anchor=base,inner sep=0pt] at (1.600,0) {$.$};
\node[anchor=base,inner sep=0pt] at (1.920,0) {$7$};
\node[anchor=base,inner sep=0pt] at (2.720,0) {$0$};
\node[anchor=base,inner sep=0pt] at (0.000,-1.1) {$5$};
\node[anchor=base,inner sep=0pt] at (0.800,-1.1) {$4$};
\node[anchor=base,inner sep=0pt] at (1.600,-1.1) {$.$};
\node[anchor=base,inner sep=0pt] at (1.920,-1.1) {$3$};
\node[anchor=base,inner sep=0pt] at (2.720,-1.1) {$6$};
\node[anchor=base,inner sep=0pt] at (0.000,-2.3) {$8$};
\node[anchor=base,inner sep=0pt] at (0.800,-2.3) {$3$};
\node[anchor=base,inner sep=0pt] at (1.600,-2.3) {$.$};
\node[anchor=base,inner sep=0pt] at (1.920,-2.3) {$0$};
\node[anchor=base,inner sep=0pt] at (2.720,-2.3) {$6$};
\node[anchor=base,inner sep=0pt] at (-.8,-1.1) {$+$};
\draw[line width=.5pt] (-1.1,-1.52)--(3.240,-1.52);
\node[anchor=base,inner sep=0pt,font=\normalsize] at (0.800,.77) {$1$};
\node[anchor=base,inner sep=0pt,font=\normalsize] at (0.000,.77) {$1$};
\end{tikzpicture}}

\MethodBlock{方法二：拆成整数部分和小数部分}{上方竖式及题目要求的验算仍保留；本法用于核对。}{\begin{aligned}[t]&\phantom{={}}28.7+54.36\\
&=28.7+54+0.36\\
&=82.7+0.36\\
&=83.06\end{aligned}}

\MethodBlock{方法三：按相同计数单位计算}{上方竖式及题目要求的验算仍保留；本法用于核对。先看成若干个百分之一，按整数加减，最后还原成小数。}{\begin{aligned}[t]&\phantom{={}}28.7+54.36\\
&=(2870+5436)\div 100\\
&=8306\div 100\\
&=83.06\end{aligned}}
\Problem{（4）$45.2-18.78=26.42$}

\MethodTitle{方法一：按题意列竖式}
\Answer{\begin{tikzpicture}[x=1em,y=1em,baseline=(current bounding box.center),draw=answerred,text=answerred]
\node[anchor=base,inner sep=0pt] at (0.000,0) {$4$};
\node[anchor=base,inner sep=0pt] at (0.800,0) {$5$};
\node[anchor=base,inner sep=0pt] at (1.600,0) {$.$};
\node[anchor=base,inner sep=0pt] at (1.920,0) {$2$};
\node[anchor=base,inner sep=0pt] at (2.720,0) {$0$};
\node[anchor=base,inner sep=0pt] at (0.000,-1.1) {$1$};
\node[anchor=base,inner sep=0pt] at (0.800,-1.1) {$8$};
\node[anchor=base,inner sep=0pt] at (1.600,-1.1) {$.$};
\node[anchor=base,inner sep=0pt] at (1.920,-1.1) {$7$};
\node[anchor=base,inner sep=0pt] at (2.720,-1.1) {$8$};
\node[anchor=base,inner sep=0pt] at (0.000,-2.3) {$2$};
\node[anchor=base,inner sep=0pt] at (0.800,-2.3) {$6$};
\node[anchor=base,inner sep=0pt] at (1.600,-2.3) {$.$};
\node[anchor=base,inner sep=0pt] at (1.920,-2.3) {$4$};
\node[anchor=base,inner sep=0pt] at (2.720,-2.3) {$2$};
\node[anchor=base,inner sep=0pt] at (-.8,-1.1) {$-$};
\draw[line width=.5pt] (-1.1,-1.52)--(3.240,-1.52);
\draw[line width=.35pt] (-0.190,.03)--(0.190,.47);
\node[anchor=base,inner sep=0pt,font=\normalsize] at (0.000,.77) {$3$};
\draw[line width=.35pt] (0.610,.03)--(0.990,.47);
\node[anchor=base,inner sep=0pt,font=\normalsize] at (0.800,.77) {$14$};
\draw[line width=.35pt] (1.730,.03)--(2.110,.47);
\node[anchor=base,inner sep=0pt,font=\normalsize] at (1.920,.77) {$11$};
\draw[line width=.35pt] (2.530,.03)--(2.910,.47);
\node[anchor=base,inner sep=0pt,font=\normalsize] at (2.720,.77) {$10$};
\end{tikzpicture}}

\MethodBlock{方法二：多减再加回}{上方竖式及题目要求的验算仍保留；本法用于核对。减去凑成的整数后，把多减的部分加回来。}{\begin{aligned}[t]&\phantom{={}}45.2-18.78\\
&=45.2-19+0.22\\
&=26.2+0.22\\
&=26.42\end{aligned}}

\MethodBlock{方法三：拆成整数部分和小数部分}{上方竖式及题目要求的验算仍保留；本法用于核对。}{\begin{aligned}[t]&\phantom{={}}45.2-18.78\\
&=45.2-18-0.78\\
&=27.2-0.78\\
&=26.42\end{aligned}}
\Problem{（5）$6.58+13.4=19.98$}

\MethodTitle{方法一：按题意列竖式}
\Answer{\begin{tikzpicture}[x=1em,y=1em,baseline=(current bounding box.center),draw=answerred,text=answerred]
\node[anchor=base,inner sep=0pt] at (0.800,0) {$6$};
\node[anchor=base,inner sep=0pt] at (1.600,0) {$.$};
\node[anchor=base,inner sep=0pt] at (1.920,0) {$5$};
\node[anchor=base,inner sep=0pt] at (2.720,0) {$8$};
\node[anchor=base,inner sep=0pt] at (0.000,-1.1) {$1$};
\node[anchor=base,inner sep=0pt] at (0.800,-1.1) {$3$};
\node[anchor=base,inner sep=0pt] at (1.600,-1.1) {$.$};
\node[anchor=base,inner sep=0pt] at (1.920,-1.1) {$4$};
\node[anchor=base,inner sep=0pt] at (2.720,-1.1) {$0$};
\node[anchor=base,inner sep=0pt] at (0.000,-2.3) {$1$};
\node[anchor=base,inner sep=0pt] at (0.800,-2.3) {$9$};
\node[anchor=base,inner sep=0pt] at (1.600,-2.3) {$.$};
\node[anchor=base,inner sep=0pt] at (1.920,-2.3) {$9$};
\node[anchor=base,inner sep=0pt] at (2.720,-2.3) {$8$};
\node[anchor=base,inner sep=0pt] at (-.8,-1.1) {$+$};
\draw[line width=.5pt] (-1.1,-1.52)--(3.240,-1.52);
\end{tikzpicture}}

\MethodBlock{方法二：拆成整数部分和小数部分}{上方竖式及题目要求的验算仍保留；本法用于核对。}{\begin{aligned}[t]&\phantom{={}}6.58+13.4\\
&=6.58+13+0.4\\
&=19.58+0.4\\
&=19.98\end{aligned}}

\MethodBlock{方法三：按相同计数单位计算}{上方竖式及题目要求的验算仍保留；本法用于核对。先看成若干个百分之一，按整数加减，最后还原成小数。}{\begin{aligned}[t]&\phantom{={}}6.58+13.4\\
&=(658+1340)\div 100\\
&=1998\div 100\\
&=19.98\end{aligned}}
\Problem{（6）$70.25-23.4=46.85$}

\MethodTitle{方法一：按题意列竖式}
\Answer{\begin{tikzpicture}[x=1em,y=1em,baseline=(current bounding box.center),draw=answerred,text=answerred]
\node[anchor=base,inner sep=0pt] at (0.000,0) {$7$};
\node[anchor=base,inner sep=0pt] at (0.800,0) {$0$};
\node[anchor=base,inner sep=0pt] at (1.600,0) {$.$};
\node[anchor=base,inner sep=0pt] at (1.920,0) {$2$};
\node[anchor=base,inner sep=0pt] at (2.720,0) {$5$};
\node[anchor=base,inner sep=0pt] at (0.000,-1.1) {$2$};
\node[anchor=base,inner sep=0pt] at (0.800,-1.1) {$3$};
\node[anchor=base,inner sep=0pt] at (1.600,-1.1) {$.$};
\node[anchor=base,inner sep=0pt] at (1.920,-1.1) {$4$};
\node[anchor=base,inner sep=0pt] at (2.720,-1.1) {$0$};
\node[anchor=base,inner sep=0pt] at (0.000,-2.3) {$4$};
\node[anchor=base,inner sep=0pt] at (0.800,-2.3) {$6$};
\node[anchor=base,inner sep=0pt] at (1.600,-2.3) {$.$};
\node[anchor=base,inner sep=0pt] at (1.920,-2.3) {$8$};
\node[anchor=base,inner sep=0pt] at (2.720,-2.3) {$5$};
\node[anchor=base,inner sep=0pt] at (-.8,-1.1) {$-$};
\draw[line width=.5pt] (-1.1,-1.52)--(3.240,-1.52);
\draw[line width=.35pt] (-0.190,.03)--(0.190,.47);
\node[anchor=base,inner sep=0pt,font=\normalsize] at (0.000,.77) {$6$};
\draw[line width=.35pt] (0.610,.03)--(0.990,.47);
\node[anchor=base,inner sep=0pt,font=\normalsize] at (0.800,.77) {$9$};
\draw[line width=.35pt] (1.730,.03)--(2.110,.47);
\node[anchor=base,inner sep=0pt,font=\normalsize] at (1.920,.77) {$12$};
\end{tikzpicture}}

\MethodBlock{方法二：拆成整数部分和小数部分}{上方竖式及题目要求的验算仍保留；本法用于核对。}{\begin{aligned}[t]&\phantom{={}}70.25-23.4\\
&=70.25-23-0.4\\
&=47.25-0.4\\
&=46.85\end{aligned}}

\MethodBlock{方法三：按相同计数单位计算}{上方竖式及题目要求的验算仍保留；本法用于核对。先看成若干个百分之一，按整数加减，最后还原成小数。}{\begin{aligned}[t]&\phantom{={}}70.25-23.4\\
&=(7025-2340)\div 100\\
&=4685\div 100\\
&=46.85\end{aligned}}
\par\noindent

\Answer{

\QHeading{计算过程：第2题：列竖式计算的完整过程}
小数点对齐，缺位补零；从右往左逐列计算。每次进位、退位都在下面列出。
\par\medskip\noindent
\textbf{第1小题：\(35.68+14.32=50\)}\par\smallskip
\noindent
\begin{tikzpicture}[x=1em,y=1em,baseline=(current bounding box.center),draw=answerred,text=answerred]
\node[anchor=base,inner sep=0pt] at (0.000,0) {$3$};
\node[anchor=base,inner sep=0pt] at (0.800,0) {$5$};
\node[anchor=base,inner sep=0pt] at (1.600,0) {$.$};
\node[anchor=base,inner sep=0pt] at (1.900,0) {$6$};
\node[anchor=base,inner sep=0pt] at (2.700,0) {$8$};
\node[anchor=base,inner sep=0pt] at (0.000,-1.1) {$1$};
\node[anchor=base,inner sep=0pt] at (0.800,-1.1) {$4$};
\node[anchor=base,inner sep=0pt] at (1.600,-1.1) {$.$};
\node[anchor=base,inner sep=0pt] at (1.900,-1.1) {$3$};
\node[anchor=base,inner sep=0pt] at (2.700,-1.1) {$2$};
\node[anchor=base,inner sep=0pt] at (0.000,-2.3) {$5$};
\node[anchor=base,inner sep=0pt] at (0.800,-2.3) {$0$};
\node[anchor=base,inner sep=0pt] at (1.600,-2.3) {$.$};
\node[anchor=base,inner sep=0pt] at (1.900,-2.3) {$0$};
\node[anchor=base,inner sep=0pt] at (2.700,-2.3) {$0$};
\node[anchor=base,inner sep=0pt] at (-.7,-1.1) {$+$};
\draw[line width=.5pt] (-1,-1.52)--(3.300,-1.52);
\node[anchor=base,inner sep=0pt,font=\normalsize] at (1.900,.72) {$1$};
\node[anchor=base,inner sep=0pt,font=\normalsize] at (0.800,.72) {$1$};
\node[anchor=base,inner sep=0pt,font=\normalsize] at (0.000,.72) {$1$};
\end{tikzpicture}
\par\smallskip{\normalsize 上方小字表示从后一位进来的数。\par}
\vspace{0pt}
\begin{enumerate}[leftmargin=1.7em,itemsep=5pt,topsep=0pt]
\item 百分位：\(\begin{aligned}[t]&\phantom{={}}8+2\\&=10\end{aligned}\)\quad 本位写0，向前一位进1。
\item 十分位：\(\begin{aligned}[t]&\phantom{={}}6+3+1\\&=9+1\\&=10\end{aligned}\)\quad 本位写0，向前一位进1。
\item 个位：\(\begin{aligned}[t]&\phantom{={}}5+4+1\\&=9+1\\&=10\end{aligned}\)\quad 本位写0，向前一位进1。
\item 十位：\(\begin{aligned}[t]&\phantom{={}}3+1+1\\&=4+1\\&=5\end{aligned}\)\quad 本位写5。
\end{enumerate}
\par\smallskip
因此，\(\begin{aligned}[t]&\phantom{={}}35.68+14.32\\&=50.00\end{aligned}\)。
\par\medskip

\par\medskip\noindent
\textbf{第2小题：\(100.00-46.75=53.25\)}\par\smallskip
\noindent
\begin{tikzpicture}[x=1em,y=1em,baseline=(current bounding box.center),draw=answerred,text=answerred]
\node[anchor=base,inner sep=0pt] at (0.000,0) {$1$};
\node[anchor=base,inner sep=0pt] at (0.800,0) {$0$};
\node[anchor=base,inner sep=0pt] at (1.600,0) {$0$};
\node[anchor=base,inner sep=0pt] at (2.400,0) {$.$};
\node[anchor=base,inner sep=0pt] at (2.700,0) {$0$};
\node[anchor=base,inner sep=0pt] at (3.500,0) {$0$};
\node[anchor=base,inner sep=0pt] at (0.800,-1.1) {$4$};
\node[anchor=base,inner sep=0pt] at (1.600,-1.1) {$6$};
\node[anchor=base,inner sep=0pt] at (2.400,-1.1) {$.$};
\node[anchor=base,inner sep=0pt] at (2.700,-1.1) {$7$};
\node[anchor=base,inner sep=0pt] at (3.500,-1.1) {$5$};
\node[anchor=base,inner sep=0pt] at (0.800,-2.3) {$5$};
\node[anchor=base,inner sep=0pt] at (1.600,-2.3) {$3$};
\node[anchor=base,inner sep=0pt] at (2.400,-2.3) {$.$};
\node[anchor=base,inner sep=0pt] at (2.700,-2.3) {$2$};
\node[anchor=base,inner sep=0pt] at (3.500,-2.3) {$5$};
\node[anchor=base,inner sep=0pt] at (-.7,-1.1) {$-$};
\draw[line width=.5pt] (-1,-1.52)--(4.100,-1.52);
\draw[line width=.35pt] (-0.160,.03)--(0.160,.43);
\node[anchor=base,inner sep=0pt,font=\normalsize] at (0.000,.72) {$0$};
\draw[line width=.35pt] (0.640,.03)--(0.960,.43);
\node[anchor=base,inner sep=0pt,font=\normalsize] at (0.800,.72) {$9$};
\draw[line width=.35pt] (1.440,.03)--(1.760,.43);
\node[anchor=base,inner sep=0pt,font=\normalsize] at (1.600,.72) {$9$};
\draw[line width=.35pt] (2.540,.03)--(2.860,.43);
\node[anchor=base,inner sep=0pt,font=\normalsize] at (2.700,.72) {$9$};
\draw[line width=.35pt] (3.340,.03)--(3.660,.43);
\node[anchor=base,inner sep=0pt,font=\normalsize] at (3.500,.72) {$10$};
\end{tikzpicture}
\par\smallskip{\normalsize 划线数字是退位前的数，上方小字是退位后参与计算的数。\par}
\vspace{0pt}
\begin{enumerate}[leftmargin=1.7em,itemsep=5pt,topsep=0pt]
\item 百位退位：\(\begin{aligned}[t]&\phantom{={}}1-1\\&=0\end{aligned}\)\quad 借出1给后一位。
\item 十位传递退位：\(\begin{aligned}[t]&\phantom{={}}0+10-1\\&=10-1\\&=9\end{aligned}\)\quad 收到10，再借出1给后一位。
\item 个位传递退位：\(\begin{aligned}[t]&\phantom{={}}0+10-1\\&=10-1\\&=9\end{aligned}\)\quad 收到10，再借出1给后一位。
\item 十分位传递退位：\(\begin{aligned}[t]&\phantom{={}}0+10-1\\&=10-1\\&=9\end{aligned}\)\quad 收到10，再借出1给后一位。
\item 百分位接收退位：\(\begin{aligned}[t]&\phantom{={}}0+10\\&=10\end{aligned}\)\quad 前一位的1变成本位的10。
\item 百分位：\(\begin{aligned}[t]&\phantom{={}}10-5\\&=5\end{aligned}\)\quad 本位写5。
\item 十分位：\(\begin{aligned}[t]&\phantom{={}}9-7\\&=2\end{aligned}\)\quad 本位写2。
\item 个位：\(\begin{aligned}[t]&\phantom{={}}9-6\\&=3\end{aligned}\)\quad 本位写3。
\item 十位：\(\begin{aligned}[t]&\phantom{={}}9-4\\&=5\end{aligned}\)\quad 本位写5。
\item 百位：\(\begin{aligned}[t]&\phantom{={}}0-0\\&=0\end{aligned}\)\quad 本位写0。
\end{enumerate}
\par\smallskip
因此，\(\begin{aligned}[t]&\phantom{={}}100.00-46.75\\&=53.25\end{aligned}\)。
\par\medskip

\par\medskip\noindent
\textbf{第3小题：\(28.70+54.36=83.06\)}\par\smallskip
\noindent
\begin{tikzpicture}[x=1em,y=1em,baseline=(current bounding box.center),draw=answerred,text=answerred]
\node[anchor=base,inner sep=0pt] at (0.000,0) {$2$};
\node[anchor=base,inner sep=0pt] at (0.800,0) {$8$};
\node[anchor=base,inner sep=0pt] at (1.600,0) {$.$};
\node[anchor=base,inner sep=0pt] at (1.900,0) {$7$};
\node[anchor=base,inner sep=0pt] at (2.700,0) {$0$};
\node[anchor=base,inner sep=0pt] at (0.000,-1.1) {$5$};
\node[anchor=base,inner sep=0pt] at (0.800,-1.1) {$4$};
\node[anchor=base,inner sep=0pt] at (1.600,-1.1) {$.$};
\node[anchor=base,inner sep=0pt] at (1.900,-1.1) {$3$};
\node[anchor=base,inner sep=0pt] at (2.700,-1.1) {$6$};
\node[anchor=base,inner sep=0pt] at (0.000,-2.3) {$8$};
\node[anchor=base,inner sep=0pt] at (0.800,-2.3) {$3$};
\node[anchor=base,inner sep=0pt] at (1.600,-2.3) {$.$};
\node[anchor=base,inner sep=0pt] at (1.900,-2.3) {$0$};
\node[anchor=base,inner sep=0pt] at (2.700,-2.3) {$6$};
\node[anchor=base,inner sep=0pt] at (-.7,-1.1) {$+$};
\draw[line width=.5pt] (-1,-1.52)--(3.300,-1.52);
\node[anchor=base,inner sep=0pt,font=\normalsize] at (0.800,.72) {$1$};
\node[anchor=base,inner sep=0pt,font=\normalsize] at (0.000,.72) {$1$};
\end{tikzpicture}
\par\smallskip{\normalsize 上方小字表示从后一位进来的数。\par}
\vspace{0pt}
\begin{enumerate}[leftmargin=1.7em,itemsep=5pt,topsep=0pt]
\item 百分位：\(\begin{aligned}[t]&\phantom{={}}0+6\\&=6\end{aligned}\)\quad 本位写6。
\item 十分位：\(\begin{aligned}[t]&\phantom{={}}7+3\\&=10\end{aligned}\)\quad 本位写0，向前一位进1。
\item 个位：\(\begin{aligned}[t]&\phantom{={}}8+4+1\\&=12+1\\&=13\end{aligned}\)\quad 本位写3，向前一位进1。
\item 十位：\(\begin{aligned}[t]&\phantom{={}}2+5+1\\&=7+1\\&=8\end{aligned}\)\quad 本位写8。
\end{enumerate}
\par\smallskip
因此，\(\begin{aligned}[t]&\phantom{={}}28.70+54.36\\&=83.06\end{aligned}\)。
\par\medskip

\par\medskip\noindent
\textbf{第4小题：\(45.20-18.78=26.42\)}\par\smallskip
\noindent
\begin{tikzpicture}[x=1em,y=1em,baseline=(current bounding box.center),draw=answerred,text=answerred]
\node[anchor=base,inner sep=0pt] at (0.000,0) {$4$};
\node[anchor=base,inner sep=0pt] at (0.800,0) {$5$};
\node[anchor=base,inner sep=0pt] at (1.600,0) {$.$};
\node[anchor=base,inner sep=0pt] at (1.900,0) {$2$};
\node[anchor=base,inner sep=0pt] at (2.700,0) {$0$};
\node[anchor=base,inner sep=0pt] at (0.000,-1.1) {$1$};
\node[anchor=base,inner sep=0pt] at (0.800,-1.1) {$8$};
\node[anchor=base,inner sep=0pt] at (1.600,-1.1) {$.$};
\node[anchor=base,inner sep=0pt] at (1.900,-1.1) {$7$};
\node[anchor=base,inner sep=0pt] at (2.700,-1.1) {$8$};
\node[anchor=base,inner sep=0pt] at (0.000,-2.3) {$2$};
\node[anchor=base,inner sep=0pt] at (0.800,-2.3) {$6$};
\node[anchor=base,inner sep=0pt] at (1.600,-2.3) {$.$};
\node[anchor=base,inner sep=0pt] at (1.900,-2.3) {$4$};
\node[anchor=base,inner sep=0pt] at (2.700,-2.3) {$2$};
\node[anchor=base,inner sep=0pt] at (-.7,-1.1) {$-$};
\draw[line width=.5pt] (-1,-1.52)--(3.300,-1.52);
\draw[line width=.35pt] (-0.160,.03)--(0.160,.43);
\node[anchor=base,inner sep=0pt,font=\normalsize] at (0.000,.72) {$3$};
\draw[line width=.35pt] (0.640,.03)--(0.960,.43);
\node[anchor=base,inner sep=0pt,font=\normalsize] at (0.800,.72) {$14$};
\draw[line width=.35pt] (1.740,.03)--(2.060,.43);
\node[anchor=base,inner sep=0pt,font=\normalsize] at (1.900,.72) {$11$};
\draw[line width=.35pt] (2.540,.03)--(2.860,.43);
\node[anchor=base,inner sep=0pt,font=\normalsize] at (2.700,.72) {$10$};
\end{tikzpicture}
\par\smallskip{\normalsize 划线数字是退位前的数，上方小字是退位后参与计算的数。\par}
\vspace{0pt}
\begin{enumerate}[leftmargin=1.7em,itemsep=5pt,topsep=0pt]
\item 十分位退位：\(\begin{aligned}[t]&\phantom{={}}2-1\\&=1\end{aligned}\)\quad 借出1给后一位。
\item 百分位接收退位：\(\begin{aligned}[t]&\phantom{={}}0+10\\&=10\end{aligned}\)\quad 前一位的1变成本位的10。
\item 百分位：\(\begin{aligned}[t]&\phantom{={}}10-8\\&=2\end{aligned}\)\quad 本位写2。
\item 个位退位：\(\begin{aligned}[t]&\phantom{={}}5-1\\&=4\end{aligned}\)\quad 借出1给后一位。
\item 十分位接收退位：\(\begin{aligned}[t]&\phantom{={}}1+10\\&=11\end{aligned}\)\quad 前一位的1变成本位的10。
\item 十分位：\(\begin{aligned}[t]&\phantom{={}}11-7\\&=4\end{aligned}\)\quad 本位写4。
\item 十位退位：\(\begin{aligned}[t]&\phantom{={}}4-1\\&=3\end{aligned}\)\quad 借出1给后一位。
\item 个位接收退位：\(\begin{aligned}[t]&\phantom{={}}4+10\\&=14\end{aligned}\)\quad 前一位的1变成本位的10。
\item 个位：\(\begin{aligned}[t]&\phantom{={}}14-8\\&=6\end{aligned}\)\quad 本位写6。
\item 十位：\(\begin{aligned}[t]&\phantom{={}}3-1\\&=2\end{aligned}\)\quad 本位写2。
\end{enumerate}
\par\smallskip
因此，\(\begin{aligned}[t]&\phantom{={}}45.20-18.78\\&=26.42\end{aligned}\)。
\par\medskip

\par\medskip\noindent
\textbf{第5小题：\(6.58+13.40=19.98\)}\par\smallskip
\noindent
\begin{tikzpicture}[x=1em,y=1em,baseline=(current bounding box.center),draw=answerred,text=answerred]
\node[anchor=base,inner sep=0pt] at (0.800,0) {$6$};
\node[anchor=base,inner sep=0pt] at (1.600,0) {$.$};
\node[anchor=base,inner sep=0pt] at (1.900,0) {$5$};
\node[anchor=base,inner sep=0pt] at (2.700,0) {$8$};
\node[anchor=base,inner sep=0pt] at (0.000,-1.1) {$1$};
\node[anchor=base,inner sep=0pt] at (0.800,-1.1) {$3$};
\node[anchor=base,inner sep=0pt] at (1.600,-1.1) {$.$};
\node[anchor=base,inner sep=0pt] at (1.900,-1.1) {$4$};
\node[anchor=base,inner sep=0pt] at (2.700,-1.1) {$0$};
\node[anchor=base,inner sep=0pt] at (0.000,-2.3) {$1$};
\node[anchor=base,inner sep=0pt] at (0.800,-2.3) {$9$};
\node[anchor=base,inner sep=0pt] at (1.600,-2.3) {$.$};
\node[anchor=base,inner sep=0pt] at (1.900,-2.3) {$9$};
\node[anchor=base,inner sep=0pt] at (2.700,-2.3) {$8$};
\node[anchor=base,inner sep=0pt] at (-.7,-1.1) {$+$};
\draw[line width=.5pt] (-1,-1.52)--(3.300,-1.52);
\end{tikzpicture}
\par\smallskip{\normalsize 上方小字表示从后一位进来的数。\par}
\vspace{0pt}
\begin{enumerate}[leftmargin=1.7em,itemsep=5pt,topsep=0pt]
\item 百分位：\(\begin{aligned}[t]&\phantom{={}}8+0\\&=8\end{aligned}\)\quad 本位写8。
\item 十分位：\(\begin{aligned}[t]&\phantom{={}}5+4\\&=9\end{aligned}\)\quad 本位写9。
\item 个位：\(\begin{aligned}[t]&\phantom{={}}6+3\\&=9\end{aligned}\)\quad 本位写9。
\item 十位：\(\begin{aligned}[t]&\phantom{={}}0+1\\&=1\end{aligned}\)\quad 本位写1。
\end{enumerate}
\par\smallskip
因此，\(\begin{aligned}[t]&\phantom{={}}6.58+13.40\\&=19.98\end{aligned}\)。
\par\medskip

\par\medskip\noindent
\textbf{第6小题：\(70.25-23.40=46.85\)}\par\smallskip
\noindent
\begin{tikzpicture}[x=1em,y=1em,baseline=(current bounding box.center),draw=answerred,text=answerred]
\node[anchor=base,inner sep=0pt] at (0.000,0) {$7$};
\node[anchor=base,inner sep=0pt] at (0.800,0) {$0$};
\node[anchor=base,inner sep=0pt] at (1.600,0) {$.$};
\node[anchor=base,inner sep=0pt] at (1.900,0) {$2$};
\node[anchor=base,inner sep=0pt] at (2.700,0) {$5$};
\node[anchor=base,inner sep=0pt] at (0.000,-1.1) {$2$};
\node[anchor=base,inner sep=0pt] at (0.800,-1.1) {$3$};
\node[anchor=base,inner sep=0pt] at (1.600,-1.1) {$.$};
\node[anchor=base,inner sep=0pt] at (1.900,-1.1) {$4$};
\node[anchor=base,inner sep=0pt] at (2.700,-1.1) {$0$};
\node[anchor=base,inner sep=0pt] at (0.000,-2.3) {$4$};
\node[anchor=base,inner sep=0pt] at (0.800,-2.3) {$6$};
\node[anchor=base,inner sep=0pt] at (1.600,-2.3) {$.$};
\node[anchor=base,inner sep=0pt] at (1.900,-2.3) {$8$};
\node[anchor=base,inner sep=0pt] at (2.700,-2.3) {$5$};
\node[anchor=base,inner sep=0pt] at (-.7,-1.1) {$-$};
\draw[line width=.5pt] (-1,-1.52)--(3.300,-1.52);
\draw[line width=.35pt] (-0.160,.03)--(0.160,.43);
\node[anchor=base,inner sep=0pt,font=\normalsize] at (0.000,.72) {$6$};
\draw[line width=.35pt] (0.640,.03)--(0.960,.43);
\node[anchor=base,inner sep=0pt,font=\normalsize] at (0.800,.72) {$9$};
\draw[line width=.35pt] (1.740,.03)--(2.060,.43);
\node[anchor=base,inner sep=0pt,font=\normalsize] at (1.900,.72) {$12$};
\end{tikzpicture}
\par\smallskip{\normalsize 划线数字是退位前的数，上方小字是退位后参与计算的数。\par}
\vspace{0pt}
\begin{enumerate}[leftmargin=1.7em,itemsep=5pt,topsep=0pt]
\item 百分位：\(\begin{aligned}[t]&\phantom{={}}5-0\\&=5\end{aligned}\)\quad 本位写5。
\item 十位退位：\(\begin{aligned}[t]&\phantom{={}}7-1\\&=6\end{aligned}\)\quad 借出1给后一位。
\item 个位传递退位：\(\begin{aligned}[t]&\phantom{={}}0+10-1\\&=10-1\\&=9\end{aligned}\)\quad 收到10，再借出1给后一位。
\item 十分位接收退位：\(\begin{aligned}[t]&\phantom{={}}2+10\\&=12\end{aligned}\)\quad 前一位的1变成本位的10。
\item 十分位：\(\begin{aligned}[t]&\phantom{={}}12-4\\&=8\end{aligned}\)\quad 本位写8。
\item 个位：\(\begin{aligned}[t]&\phantom{={}}9-3\\&=6\end{aligned}\)\quad 本位写6。
\item 十位：\(\begin{aligned}[t]&\phantom{={}}6-2\\&=4\end{aligned}\)\quad 本位写4。
\end{enumerate}
\par\smallskip
因此，\(\begin{aligned}[t]&\phantom{={}}70.25-23.40\\&=46.85\end{aligned}\)。
\par\medskip

}

\QHeading{第3题　脱式计算。}

\CalcBlock{1}{\begin{aligned}[t]&\phantom{={}}\color{black}360-120\div6\times5\\&\color{answerred}=360-20\times5\\&\color{answerred}=360-100\\&\color{answerred}=260\end{aligned}}
\par

\CalcBlock[方法一]{2}{\begin{aligned}[t]&\phantom{={}}\color{black}45\times(120-80\div4)\\&\color{answerred}=45\times(120-20)\\&\color{answerred}=45\times100\\&\color{answerred}=4500\end{aligned}}
\MethodBlock{方法二：乘法分配律}{把括号中的两部分分别相乘，再相加或相减。}{\begin{aligned}[t]&\phantom{={}}45\times (120-80\div 4)\\
&=45\times 120-45\times 80\div 4\\
&=5400-3600\div 4\\
&=5400-900\\
&=4500\end{aligned}}

\CalcBlock{3}{\begin{aligned}[t]&\phantom{={}}\color{black}(240+180)\div30\times15\\&\color{answerred}=420\div30\times15\\&\color{answerred}=14\times15\\&\color{answerred}=210\end{aligned}}
\CalcBlock[方法一]{4}{\begin{aligned}[t]&\phantom{={}}\color{black}780-(260+15\times12)\\&\color{answerred}=780-(260+180)\\&\color{answerred}=780-440\\&\color{answerred}=340\end{aligned}}
\MethodBlock{方法二：先处理括号或乘除，再交换、结合后先凑整}{保留每个数原来的加减号，先把能凑整的部分放在一起。}{\begin{aligned}[t]&\phantom{={}}780-(260+15\times 12)\\
&=780-(260+180)\\
&=(780-180)-260\\
&=600-260\\
&=340\end{aligned}}

\MethodBlock{方法三：先处理括号或乘除，再去括号后计算}{加号后的括号去掉不变号，减号后的括号去掉要改变各项的加减号。}{\begin{aligned}[t]&\phantom{={}}780-(260+15\times 12)\\
&=780-(260+180)\\
&=780-260-180\\
&=520-180\\
&=340\end{aligned}}

\CalcBlock{5}{\begin{aligned}[t]&\phantom{={}}\color{black}450+150\div30\times5\\&\color{answerred}=450+5\times5\\&\color{answerred}=450+25\\&\color{answerred}=475\end{aligned}}
\CalcBlock[方法一]{6}{\begin{aligned}[t]&\phantom{={}}\color{black}(420-180)\div(12\times5)\\&\color{answerred}=240\div60\\&\color{answerred}=4\end{aligned}}
\MethodBlock{方法二：除以积，改成连续除}{连续除以两个数，等于除以这两个数的积。}{\begin{aligned}[t]&\phantom{={}}(420-180)\div (12\times 5)\\
&=(420-180)\div 12\div 5\\
&=240\div 12\div 5\\
&=20\div 5\\
&=4\end{aligned}}

\MethodBlock{方法三：先处理括号或乘除，再先分别相除再合并}{}{\begin{aligned}[t]&\phantom{={}}(420-180)\div (12\times 5)\\
&=(420-180)\div 60\\
&=420\div 60-180\div 60\\
&=7-3\\
&=4\end{aligned}}


